\chapter{Trigonométrie~: le cercle trigonométrique}\label{ch:g11:trigo}

La trigonométrie commence dans les triangles rectangles, mais sa vraie
demeure est le \emph{cercle trigonométrique}, où le \omterm{def:g11:trigo:cossin}{cosinus} et le \omterm{def:g11:trigo:cossin}{sinus}
deviennent des \omterm{def:g11:func:function}{fonctions} d'un \omterm{def:g10:numbers:sets}{nombre réel} arbitraire. Ce chapitre met en
place le \omterm{def:g11:trigo:radian}{radian}, l'enroulement de la droite réelle autour du cercle, les
valeurs exactes à connaître par cœur, et les symétries qui engendrent
toutes les identités classiques. L'étude de $\cos$ et $\sin$ comme
\omterm{def:g11:func:function}{fonctions} --- variations, \omterm{def:g11:deriv:derivative}{dérivées} --- est menée dans le
\cref{ch:g12:trigo}.

\section{Radians et enroulement}

\begin{definition}[Radian]\label{def:g11:trigo:radian}
Soit $\mathcal C$ le cercle de rayon $1$ centré à l'origine $O$, et
$I(1, 0)$. La mesure en \emph{radians}\index{radian} d'un angle au centre
de $\mathcal C$ est la longueur de l'arc qu'il découpe sur $\mathcal C$.
Comme le cercle complet a pour circonférence $2\pi$~:
\[
360^\circ = 2\pi \text{ rad},
\qquad
180^\circ = \pi \text{ rad},
\qquad
1 \text{ rad} = \frac{180^\circ}{\pi} \approx 57.3^\circ .
\]
\end{definition}

\begin{method}[Convertir]\label{met:g11:trigo:convert}
Degrés et \omterm{def:g11:trigo:radian}{radians} sont proportionnels~: multiplier par $\frac{\pi}{180}$
pour passer des degrés aux \omterm{def:g11:trigo:radian}{radians}, par $\frac{180}{\pi}$ dans l'autre
sens. Par exemple $60^\circ = 60 \times \frac{\pi}{180} = \frac{\pi}{3}$,
et $\frac{3\pi}{4} = \frac{3 \times 180^\circ}{4} = 135^\circ$.
\end{method}

\begin{definition}[Enrouler la droite sur le cercle]\label{def:g11:trigo:winding}
À chaque \omterm{def:g10:numbers:sets}{nombre réel} $t$, associer le point $M(t)$ de $\mathcal C$
obtenu en parcourant une distance $\abs t$ le long du cercle depuis $I$,
dans le sens trigonométrique si $t \geq 0$, dans le sens horaire si
$t < 0$. Comme la circonférence est $2\pi$, les nombres $t$ et
$t + 2k\pi$ ($k \in \Z$) atteignent le même point~:
$M(t + 2k\pi) = M(t)$.
\end{definition}

\begin{example}
$M(0) = I$~; $M\!\left(\frac{\pi}{2}\right) = (0, 1)$, le sommet du
cercle~; $M(\pi) = (-1, 0)$~; $M\!\left(-\frac{\pi}{2}\right) =
(0, -1)$~; et $M\!\left(\frac{9\pi}{4}\right) =
M\!\left(\frac{\pi}{4} + 2\pi\right) = M\!\left(\frac{\pi}{4}\right)$.
\end{example}

\section{Cosinus et sinus}

\begin{definition}[Cosinus et sinus]\label{def:g11:trigo:cossin}
Pour tout réel $t$, le \emph{cosinus}\index{cosinus} et le
\emph{sinus}\index{sinus} de $t$ sont les \omterm{def:g10:coordgeom:system}{coordonnées} du point $M(t)$ du
cercle trigonométrique~:
\[
M(t) = (\cos t,\ \sin t).
\]
\end{definition}

\begin{omfigure}
\begin{tikzpicture}[scale=2.0]
  \draw[->] (-1.3,0) -- (1.4,0) node[right] {$x$};
  \draw[->] (0,-1.3) -- (0,1.4) node[above] {$y$};
  \draw[gray] (0,0) circle (1);
  \coordinate (O) at (0,0);
  \coordinate (M) at (55:1);
  \draw[thick,omDef] (O) -- (M);
  \draw[dashed] (M) -- (M |- O) node[below, font=\small] {$\cos t$};
  \draw[dashed] (M) -- (M -| O) node[left, font=\small] {$\sin t$};
  \draw[->,omThm,thick] (0.35,0) arc (0:55:0.35);
  \node[omThm, font=\small] at (27:0.55) {$t$};
  \fill (M) circle (0.7pt) node[above right] {$M(t)$};
  \fill (1,0) circle (0.7pt) node[below right] {$I$};
\end{tikzpicture}

{\small Le cercle trigonométrique~: le \omterm{def:g10:numbers:sets}{nombre réel} $t$, enroulé dans le
sens trigonométrique depuis $I$, arrive au point $M(t)$ dont les
\omterm{def:g10:coordgeom:system}{coordonnées} \emph{définissent} $\cos t$ et $\sin t$.}
\end{omfigure}

\begin{proposition}[Premières propriétés]\label{prop:g11:trigo:first}
Pour tous $t \in \R$ et $k \in \Z$~:
\[
-1 \leq \cos t \leq 1, \qquad
-1 \leq \sin t \leq 1, \qquad
\cos^2 t + \sin^2 t = 1,
\]
\[
\cos(t + 2k\pi) = \cos t, \qquad \sin(t + 2k\pi) = \sin t .
\]
\end{proposition}

\begin{proof}
$M(t)$ est sur le cercle unité, donc ses \omterm{def:g10:coordgeom:system}{coordonnées} sont entre $-1$ et
$1$, et elles satisfont l'\omterm{def:g10:algebra:equation}{équation} du cercle $x^2 + y^2 = 1$, qui est
l'identité $\cos^2 t + \sin^2 t = 1$. La périodicité reformule
$M(t + 2k\pi) = M(t)$ (\cref{def:g11:trigo:winding}).
\end{proof}

\begin{example}\label{ex:g11:trigo:identity}
Si $\sin t = \frac35$ et $t \in \intcc{\frac\pi2}{\pi}$ (deuxième
quadrant), alors $\cos^2 t = 1 - \frac{9}{25} = \frac{16}{25}$, donc
$\cos t = \pm\frac45$~; dans le deuxième quadrant l'\omterm{def:g10:coordgeom:system}{abscisse} est
négative, d'où $\cos t = -\frac45$.
\end{example}

Les valeurs à connaître par cœur~:
\begin{center}
\begin{tabular}{c|ccccc}
$t$ & $0$ & $\dfrac{\pi}{6}$ & $\dfrac{\pi}{4}$ & $\dfrac{\pi}{3}$ & $\dfrac{\pi}{2}$\\[6pt]
\hline\rule{0pt}{14pt}%
$\cos t$ & $1$ & $\dfrac{\sqrt3}{2}$ & $\dfrac{\sqrt2}{2}$ & $\dfrac12$ & $0$\\[6pt]
$\sin t$ & $0$ & $\dfrac12$ & $\dfrac{\sqrt2}{2}$ & $\dfrac{\sqrt3}{2}$ & $1$
\end{tabular}
\end{center}

\begin{remark}
Un moyen mnémotechnique~: les \omterm{def:g11:trigo:cossin}{sinus} se lisent
$\frac{\sqrt0}{2}, \frac{\sqrt1}{2}, \frac{\sqrt2}{2}, \frac{\sqrt3}{2},
\frac{\sqrt4}{2}$, et les \omterm{def:g11:trigo:cossin}{cosinus} sont la même liste inversée.
\end{remark}

\section{Angles associés}

\begin{proposition}[Angles associés]\label{prop:g11:trigo:associated}
\index{angles associés}
Pour tout $t \in \R$~:
\[
\begin{aligned}
\cos(-t) &= \cos t, & \sin(-t) &= -\sin t,\\
\cos(\pi - t) &= -\cos t, & \sin(\pi - t) &= \sin t,\\
\cos(\pi + t) &= -\cos t, & \sin(\pi + t) &= -\sin t,\\
\cos\left(\tfrac{\pi}{2} - t\right) &= \sin t, &
\sin\left(\tfrac{\pi}{2} - t\right) &= \cos t .
\end{aligned}
\]
\end{proposition}

\begin{proof}
Chaque ligne exprime une symétrie du cercle.

$M(-t)$ est le symétrique de $M(t)$ par rapport à l'axe des $x$~:
l'\omterm{def:g10:coordgeom:system}{abscisse} est inchangée, l'ordonnée change de signe.

$M(\pi - t)$ est le symétrique de $M(t)$ par rapport à l'axe des $y$~:
marcher en arrière depuis $\pi$ de $t$ arrive à l'opposé
(horizontalement) de marcher en avant depuis $0$ de $t$. L'\omterm{def:g10:coordgeom:system}{abscisse}
change de signe, l'ordonnée est inchangée.

$M(\pi + t)$ est diamétralement opposé à $M(t)$ (un demi-tour)~: les
deux \omterm{def:g10:coordgeom:system}{coordonnées} changent de signe.

$M\!\left(\frac\pi2 - t\right)$ est le symétrique de $M(t)$ par rapport
à la diagonale $y = x$, qui échange les deux \omterm{def:g10:coordgeom:system}{coordonnées}.
\end{proof}

\begin{omfigure}
\begin{tikzpicture}[scale=2.0]
  \draw[->] (-1.35,0) -- (1.45,0) node[right] {$x$};
  \draw[->] (0,-1.35) -- (0,1.45) node[above] {$y$};
  \draw[gray] (0,0) circle (1);
  \fill (40:1) circle (0.7pt) node[above right, font=\small] {$M(t)$};
  \fill (-40:1) circle (0.7pt) node[below right, font=\small] {$M(-t)$};
  \fill (140:1) circle (0.7pt) node[above left, font=\small]
    {$M(\pi - t)$};
  \fill (220:1) circle (0.7pt) node[below left, font=\small]
    {$M(\pi + t)$};
  \draw[black, dashed, thin] (40:1) -- (-40:1);
  \draw[black, dashed, thin] (40:1) -- (140:1);
  \draw[black, dashed, thin] (140:1) -- (220:1);
  \draw[omDef, thin] (0,0) -- (40:1);
  \draw[omDef, thin] (0,0) -- (140:1);
  \draw[omDef, thin] (0,0) -- (220:1);
  \draw[omDef, thin] (0,0) -- (-40:1);
\end{tikzpicture}

{\small Les quatre points associés~: $M(-t)$ est le miroir de $M(t)$ par
rapport à l'axe des $x$, $M(\pi - t)$ par rapport à l'axe des $y$, et
$M(\pi + t)$ est diamétralement opposé. Toute identité de la
\cref{prop:g11:trigo:associated} se lit sur ce dessin.}
\end{omfigure}

\begin{example}
$\cos\frac{2\pi}{3} = \cos\left(\pi - \frac\pi3\right)
= -\cos\frac\pi3 = -\frac12$, et
$\sin\frac{7\pi}{6} = \sin\left(\pi + \frac\pi6\right)
= -\sin\frac\pi6 = -\frac12$. Le tableau de cinq valeurs, étendu par les
symétries, couvre tout le cercle.
\end{example}

\section{Résoudre des équations trigonométriques}

\begin{theorem}[Les équations $\cos t = \cos a$ et $\sin t = \sin a$]
\label{thm:g11:trigo:equations}
Pour des réels $t$ et $a$~:
\[
\cos t = \cos a \iff t = a + 2k\pi \ \text{ou}\ t = -a + 2k\pi
\quad (k \in \Z),
\]
\[
\sin t = \sin a \iff t = a + 2k\pi \ \text{ou}\ t = \pi - a + 2k\pi
\quad (k \in \Z).
\]
\end{theorem}

\begin{proof}
Deux points du cercle trigonométrique ont la même \omterm{def:g10:coordgeom:system}{abscisse} exactement
lorsqu'ils sont égaux ou symétriques l'un de l'autre par rapport à l'axe
des $x$~; par la \cref{prop:g11:trigo:associated}, le symétrique de
$M(a)$ est $M(-a)$. Donc $\cos t = \cos a$ signifie $M(t) = M(a)$ ou
$M(t) = M(-a)$, c'est-à-dire $t = \pm a$ à un multiple de $2\pi$ près.
De même, des ordonnées égales signifient des points égaux ou symétriques
par rapport à l'axe des $y$, et le symétrique de $M(a)$ est
$M(\pi - a)$.
\end{proof}

\begin{method}[Résoudre $\cos t = c$ sur un intervalle]\label{met:g11:trigo:solve}
\begin{enumerate}
  \item Trouver \emph{un} angle $a$ avec $\cos a = c$, d'après le tableau
        des valeurs connues.
  \item Écrire les solutions générales $t = \pm a + 2k\pi$
        (\cref{thm:g11:trigo:equations}).
  \item Choisir les entiers $k$ qui tombent dans l'\omterm{def:g10:numbers:interval}{intervalle} demandé, et
        lister les solutions.
\end{enumerate}
Procéder de même pour $\sin t = c$ avec $t = a + 2k\pi$ ou
$t = \pi - a + 2k\pi$.
\end{method}

\begin{example}\label{ex:g11:trigo:solve}
Résoudre $\cos t = \frac12$ sur $\intoc{-\pi}{\pi}$. Une solution de
$\cos a = \frac12$ est $a = \frac{\pi}{3}$. Les solutions générales sont
$t = \frac\pi3 + 2k\pi$ et $t = -\frac\pi3 + 2k\pi$. Dans
$\intoc{-\pi}{\pi}$, seul $k = 0$ contribue~:
$t \in \left\{-\frac\pi3,\ \frac\pi3\right\}$.

Résoudre $\sin t = -\frac{\sqrt2}{2}$ sur $\intco{0}{2\pi}$. Un angle est
$a = -\frac\pi4$~; les solutions sont $t = -\frac\pi4 + 2k\pi$ et
$t = \pi + \frac\pi4 + 2k\pi = \frac{5\pi}{4} + 2k\pi$. Dans
$\intco{0}{2\pi}$~: $t = \frac{7\pi}{4}$ (pour $k = 1$) et
$t = \frac{5\pi}{4}$.
\end{example}

\begin{omfigure}
\begin{tikzpicture}[scale=2.0]
  \draw[->] (-1.35,0) -- (1.45,0) node[right] {$x$};
  \draw[->] (0,-1.2) -- (0,1.3) node[above] {$y$};
  \draw[gray] (0,0) circle (1);
  \draw[omProp, thick] (0.5,-1.05) -- (0.5,1.15)
    node[above, font=\small] {$x = \tfrac12$};
  \fill (60:1) circle (0.7pt)
    node[above right, font=\small] {$M\!\left(\tfrac\pi3\right)$};
  \fill (-60:1) circle (0.7pt)
    node[below right, font=\small] {$M\!\left(-\tfrac\pi3\right)$};
  \draw[omDef, thin] (0,0) -- (60:1);
  \draw[omDef, thin] (0,0) -- (-60:1);
\end{tikzpicture}

{\small Résoudre $\cos t = \frac12$~: la droite verticale $x = \frac12$
(orange) coupe le cercle en deux points, symétriques par rapport à
l'axe des $x$ --- d'où les deux familles de solutions
$t = \pm\frac\pi3 + 2k\pi$.}
\end{omfigure}

\section{Exercices}

\begin{exercise}[$\star$]\label{exo:g11:trigo:1}
Convertir en \omterm{def:g11:trigo:radian}{radians}~: $30^\circ$, $45^\circ$, $120^\circ$, $270^\circ$.
Convertir en degrés~: $\frac{\pi}{5}$, $\frac{5\pi}{6}$,
$\frac{7\pi}{4}$, $\frac{2\pi}{9}$.
\end{exercise}

\begin{exercise}[$\star$]\label{exo:g11:trigo:2}
Placer sur le cercle trigonométrique les points $M(t)$ pour
\[
t = \frac{2\pi}{3},\quad
t = -\frac{\pi}{4},\quad
t = \frac{17\pi}{6},\quad
t = -\frac{7\pi}{2},
\]
après avoir réduit chacun à une valeur dans $\intoc{-\pi}{\pi}$ modulo
$2\pi$.
\end{exercise}

\begin{exercise}[$\star$]\label{exo:g11:trigo:3}
En utilisant le tableau et les \omterm{prop:g11:trigo:associated}{angles associés}, donner les valeurs
exactes de
\[
\cos\frac{3\pi}{4}, \quad
\sin\frac{5\pi}{6}, \quad
\cos\left(-\frac{\pi}{3}\right), \quad
\sin\frac{4\pi}{3}, \quad
\cos\frac{11\pi}{6} .
\]
\end{exercise}

\begin{exercise}[$\star$]\label{exo:g11:trigo:4}
Sachant $\cos t = \frac{5}{13}$ avec $t \in \intoo{-\frac\pi2}{0}$,
calculer $\sin t$ exactement.
\end{exercise}

\begin{exercise}[$\star\star$]\label{exo:g11:trigo:5}
Simplifier, pour tout réel $x$~:
\[
A = \cos(\pi + x) + \cos(-x) + \sin\left(\frac\pi2 - x\right) - \cos x,
\qquad
B = \sin(\pi - x) + \sin(\pi + x) + \sin(-x) .
\]
\end{exercise}

\begin{exercise}[$\star\star$]\label{exo:g11:trigo:6}
Résoudre sur $\intoc{-\pi}{\pi}$~:
\[
\cos t = \frac{\sqrt3}{2}, \qquad
\sin t = \frac12, \qquad
\cos t = -1 .
\]
\end{exercise}

\begin{exercise}[$\star\star$]\label{exo:g11:trigo:7}
Résoudre sur $\intco{0}{2\pi}$~:
\[
\sin t = -\frac{\sqrt3}{2}, \qquad
\cos t = 0, \qquad
2\sin t + 1 = 0 .
\]
\end{exercise}

\begin{exercise}[$\star\star$]\label{exo:g11:trigo:8}
Résoudre sur $\R$ l'\omterm{def:g10:algebra:equation}{équation} $\cos 2t = \cos t$. (Utiliser le
\cref{thm:g11:trigo:equations} avec $a = t$, et discuter les deux
familles de solutions.)
\end{exercise}

\begin{exercise}[$\star\star$]\label{exo:g11:trigo:9}
Montrer que pour tout réel $x$,
\[
\bigl(\cos x + \sin x\bigr)^2 + \bigl(\cos x - \sin x\bigr)^2 = 2 .
\]
\end{exercise}

\begin{exercise}[$\star\star$]\label{exo:g11:trigo:10}
Une roue de rayon $30$~cm roule sans glisser. De quel angle (en \omterm{def:g11:trigo:radian}{radians},
puis en degrés) tourne-t-elle lorsque le vélo avance de $1.5$~m~?
Quelle distance le vélo parcourt-il pendant un tour complet de la roue~?
\end{exercise}

\begin{exercise}[$\star\star\star$]\label{exo:g11:trigo:11}
Résoudre sur $\intoc{-\pi}{\pi}$ l'inéquation $\cos t \leq \frac12$.
(Croquer le cercle trigonométrique, marquer la région où l'\omterm{def:g10:coordgeom:system}{abscisse} est
au plus $\frac12$, et lire l'\omterm{def:g10:numbers:interval}{intervalle} de valeurs de $t$.)
\end{exercise}

\section{Problème~: la grande roue et la marée}

\begin{problem}[{Devoir du week-end --- le radian comme distance
roulée, et le cercle trigonométrique comme horloge maîtresse de tout
phénomène périodique}]
\label{pb:g11:trigo:1}
Une grande roue vous promène sur un cercle~; la marée promène l'eau
du port de haut en bas sur le même cercle mathématique. Tout
phénomène périodique --- roues, marées, sons, saisons --- lit son
horaire sur le cercle trigonométrique de ce chapitre
(\cref{def:g11:trigo:cossin}), et les \omterm{def:g10:algebra:equation}{équations} de la
\cref{met:g11:trigo:solve} en donnent les instants exacts. Ce
problème convertit, enroule, modélise et résout --- et règle en
chemin la question de savoir pourquoi les mathématiciens mesurent
les angles en \omterm{def:g11:trigo:radian}{radians}.

\medskip
\noindent\textbf{Partie I --- Le \omterm{def:g11:trigo:radian}{radian}, une distance roulée.}
\begin{enumerate}
  \item Convertir en \omterm{def:g11:trigo:radian}{radians}~: $30^\circ$, $135^\circ$,
        $300^\circ$~; et en degrés~: $\frac{3\pi}{4}$,
        $\frac{5\pi}{6}$ (\cref{met:g11:trigo:convert}).
  \item Tout l'intérêt du \omterm{def:g11:trigo:radian}{radian}~: un arc d'angle $t$ (en \omterm{def:g11:trigo:radian}{radians})
        sur un cercle de rayon $r$ a pour longueur exactement
        $r\,t$. Quel arc un angle de $2$ \omterm{def:g11:trigo:radian}{radians} découpe-t-il sur un
        cercle de rayon $3$ m~? (Aucun $\pi$ nulle part --- c'est
        bien là le point.)
  \item Une roue de rayon $30$ cm roule sur $1$ km sans glisser. De
        combien de \emph{\omterm{def:g11:trigo:radian}{radians}} a-t-elle tourné, et cela fait
        combien de tours complets~?
  \item Placer sur le cercle trigonométrique et évaluer
        exactement~: $\cos\frac{2\pi}{3}$,
        $\sin\left(-\frac{\pi}{4}\right)$ et $\cos\frac{19\pi}{6}$
        (réduire d'abord modulo $2\pi$~;
        \cref{prop:g11:trigo:associated}).
  \item Un angle $t \in \intoo{\frac\pi2}{\pi}$ vérifie
        $\sin t = \frac35$. À l'aide de
        $\cos^2 t + \sin^2 t = 1$ et du quadrant, déterminer
        $\cos t$ et $\tan t$.
\end{enumerate}

\medskip
\noindent\textbf{Partie II --- La grande roue.}
Une grande roue a un rayon de $60$ m, son moyeu est à $65$ m du sol,
et elle fait un tour toutes les $30$ minutes. On y monte au point le
plus bas à l'instant $t = 0$ (en minutes), de sorte que la hauteur
obéit à
\[
h(t) = 65 - 60 \cos\!\left(\frac{\pi t}{15}\right).
\]
\begin{enumerate}[resume]
  \item Justifier la formule~: vérifier l'angle parcouru après $t$
        minutes, la hauteur en $t = 0$, et expliquer le signe
        moins.
  \item Calculer la hauteur à $t = 7.5$, $t = 15$ et $t = 20$
        minutes (valeurs exactes).
  \item À quel instant est-on \emph{pour la première fois} à
        $95$ m~? (Résoudre $h(t) = 95$ avec la
        \cref{met:g11:trigo:solve}.)
  \item Pendant quel \omterm{def:g10:numbers:interval}{intervalle} de temps du tour est-on à au moins
        $95$ m de hauteur --- et pendant combien de temps au total~?
        (Comparer avec l'\cref{exo:g11:trigo:11}.)
  \item Où la hauteur varie-t-elle \emph{le plus vite}~? Comparer
        la montée pendant la première minute ($h(1) - h(0)$) avec
        celle entre les minutes $7$ et $8$, et formuler
        l'observation (l'outil qui la rend exacte --- la \omterm{def:g11:deriv:derivative}{dérivée} de
        la sinusoïde --- arrive l'an prochain).
  \item À vous de jouer les ingénieurs~: concevoir la formule de
        hauteur d'une roue de hauteur maximale $40$ m, de hauteur
        minimale $4$ m et de période $12$ minutes, l'embarquement se
        faisant en bas à $t = 0$.
\end{enumerate}

\medskip
\noindent\textbf{Partie III --- La marée et les symétries du
cercle.}
Dans un port, la hauteur d'eau en mètres, $t$ heures après minuit,
vaut
\[
d(t) = 7 + 3 \sin\!\left(\frac{\pi t}{6}\right).
\]
\begin{enumerate}[resume]
  \item Donner les hauteurs d'eau maximale et minimale, les instants
        où elles se produisent pour la première fois, et la période
        de la marée.
  \item Un cargo a besoin de $8.5$ m d'eau. Résoudre
        $d(t) \geq 8.5$ sur $\intco{0}{12}$~: pendant quelle
        fenêtre peut-il entrer, et quelle est sa durée~? Quand
        s'ouvre la fenêtre suivante~?
  \item Retrouver directement, à partir des symétries du cercle,
        deux identités d'\omterm{prop:g11:trigo:associated}{angles associés} --- $\sin(\pi - t) =
        \sin t$ (symétrie par rapport à l'axe vertical) et
        $\cos(\pi + t) = -\cos t$ (demi-tour) --- puis énoncer
        l'identité des angles \omterm{def:g10:proba:operations}{complémentaires}
        $\cos\left(\frac\pi2 - t\right) = \sin t$.
  \item Parité (vocabulaire du \cref{pb:g11:func:1})~: classer
        $\sin$, $\cos$ et $t \mapsto \sin t + t^3$ en \omterm{def:g11:func:parity}{fonctions
        paires} ou \omterm{def:g11:func:parity}{impaires}, avec une démonstration d'une ligne à
        partir du cercle.
  \item Résoudre sur $\intco{0}{2\pi}$~:
        $2\sin^2 t - \sin t - 1 = 0$. (Un trinôme du second degré
        déguisé~: le \omterm{def:g10:algebra:expand}{factoriser} comme au \cref{pb:g11:quad:1}, puis
        lire le cercle.)
\end{enumerate}

\medskip
\noindent\textbf{Partie IV --- L'angle du mathématicien.}
\begin{enumerate}[resume]
  \item Combien de degrés vaut $1$ \omterm{def:g11:trigo:radian}{radian}~? Puis un piège de
        calculatrice~: qui est le plus grand, $\sin(1)$ (mode
        \omterm{def:g11:trigo:radian}{radian}) ou $\sin(1^\circ)$ --- et à peu près dans quel
        rapport~? Morale sur le réglage des calculatrices~?
  \item Résoudre $\cos t = \sin t$ sur $\intco{0}{2\pi}$ à l'aide
        du cercle (où l'\omterm{def:g10:coordgeom:system}{abscisse} et l'ordonnée sont-elles
        égales~?).
  \item Pour les petits angles, l'arc et son \omterm{def:g11:trigo:cossin}{sinus} coïncident
        presque~: comparer $\sin(0.1)$ et $0.1$ (cinq décimales),
        et donner l'erreur relative. Cette
        «~\omterm{def:g10:numbers:approx}{approximation} des petits angles~» guide navires et
        pendules~; le théorème qui la porte
        ($\frac{\sin t}{t} \to 1$) est une vedette du chapitre sur
        les limites de l'an prochain.
  \item Pour finir --- le cercle trigonométrique comme horloge
        maîtresse, une phrase par point~: le \omterm{def:g11:trigo:radian}{radian} transforme les
        angles en nombres honnêtes (arc $=$ rayon $\times$ angle)~;
        le \omterm{def:g11:trigo:cossin}{cosinus} et le \omterm{def:g11:trigo:cossin}{sinus} sont l'ombre et la hauteur de
        l'aiguille~; les symétries du cercle sont les identités~;
        résoudre $\cos t = c$ revient à lire des instants sur
        l'horloge~; et tout phénomène périodique --- roue, marée,
        son --- n'est que cette horloge en costume.
\end{enumerate}
\end{problem}
